\documentclass[14pt]{extarticle}

\usepackage[a4paper,margin=1in]{geometry}
\usepackage{amsmath,amssymb}
\usepackage{booktabs}
\usepackage{hyperref}

\def\mathbi#1{\textbf{\em #1}}

\title{Bayesian Ridge Regression}
\author{}
\date{}

\begin{document}

\maketitle

Bayesian Ridge Regression is a Bayesian extension of linear regression and ridge regression. Instead of estimating a single optimal coefficient vector, Bayesian Ridge Regression treats the model coefficients as random variables and estimates their probability distributions from the observed data.

Like ridge regression, Bayesian Ridge Regression shrinks large coefficients toward zero, reducing overfitting and improving model stability. However, unlike ridge regression, the amount of regularization is automatically inferred from the data rather than chosen manually.

Suppose we are given a dataset consisting of $n$ observations and $p$ features. Let

\[
\mathbf{X}
=
\begin{bmatrix}
x_{11} & x_{12} & \cdots & x_{1p} \\
x_{21} & x_{22} & \cdots & x_{2p} \\
\vdots & \vdots & \ddots & \vdots \\
x_{n1} & x_{n2} & \cdots & x_{np}
\end{bmatrix}
\]

be the feature matrix and

\[
\mathbi{y}
=
\begin{bmatrix}
y_1 \\
y_2 \\
\vdots \\
y_n
\end{bmatrix}
\]

be the vector of target values.

As in linear regression, after absorbing the bias term into the feature matrix, the prediction function can be written as

\[
\hat{\mathbi{y}}
=
\mathbf{X}\mathbi{w},
\]

where $\mathbf{X}\in\mathbb{R}^{n\times p}$ is the feature matrix,
$\mathbi{w}\in\mathbb{R}^{p}$ is the coefficient vector,
and $\hat{\mathbi{y}}\in\mathbb{R}^{n}$ is the vector of predicted target values.

Bayesian Ridge Regression assumes that the target values are generated according to

\[
\mathbi{y}
=
\mathbf{X}\mathbi{w}
+
\boldsymbol{\varepsilon},
\]

where the noise vector follows a Gaussian distribution

\[
\boldsymbol{\varepsilon}
\sim
\mathcal{N}
\left(
\mathbf{0},
\alpha^{-1}\mathbf{I}
\right).
\]

Consequently, the likelihood function is

\[
p(\mathbi{y}\mid \mathbf{X},\mathbi{w},\alpha)
=
\mathcal{N}
\left(
\mathbf{X}\mathbi{w},
\alpha^{-1}\mathbf{I}
\right),
\]

where $\alpha$ is called the noise precision.

To regularize the model, Bayesian Ridge Regression places a Gaussian prior on the coefficient vector:

\[
p(\mathbi{w}\mid \lambda)
=
\mathcal{N}
\left(
\mathbf{0},
\lambda^{-1}\mathbf{I}
\right),
\]

where $\lambda$ is called the weight precision.

Using Bayes' theorem, the posterior distribution of the coefficients is

\[
p(\mathbi{w}\mid \mathbf{X},\mathbi{y},\alpha,\lambda)
=
\frac{
p(\mathbi{y}\mid \mathbf{X},\mathbi{w},\alpha)
\,p(\mathbi{w}\mid \lambda)
}{
p(\mathbi{y}\mid \mathbf{X},\alpha,\lambda)
}.
\]

Substituting the Gaussian likelihood and Gaussian prior gives

\[
p(\mathbi{w}\mid \mathbf{X},\mathbi{y},\alpha,\lambda)
\propto
\exp
\left(
-\frac{\alpha}{2}
\|\mathbf{X}\mathbi{w}-\mathbi{y}\|^2
-\frac{\lambda}{2}
\|\mathbi{w}\|^2
\right).
\]

The posterior distribution is also Gaussian:

\[
p(\mathbi{w}\mid \mathbf{X},\mathbi{y},\alpha,\lambda)
=
\mathcal{N}
(
\boldsymbol{\mu},
\mathbf{\Sigma}
).
\]

The posterior covariance matrix is

\[
\mathbf{\Sigma}
=
\left(
\lambda\mathbf{I}
+
\alpha\mathbf{X}^{\top}\mathbf{X}
\right)^{-1},
\]

and the posterior mean is

\[
\boldsymbol{\mu}
=
\alpha
\mathbf{\Sigma}
\mathbf{X}^{\top}
\mathbi{y}.
\]

Substituting the expression for $\mathbf{\Sigma}$ yields

\[
\boldsymbol{\mu}
=
\left(
\mathbf{X}^{\top}\mathbf{X}
+
\frac{\lambda}{\alpha}\mathbf{I}
\right)^{-1}
\mathbf{X}^{\top}
\mathbi{y}.
\]

Notice that this expression has the same form as the ridge regression solution

\[
\mathbi{w}
=
\left(
\mathbf{X}^{\top}\mathbf{X}
+
\gamma\mathbf{I}
\right)^{-1}
\mathbf{X}^{\top}
\mathbi{y},
\]

where

\[
\gamma
=
\frac{\lambda}{\alpha}.
\]

Therefore, ridge regression can be interpreted as finding the maximum a posteriori estimate of the coefficients under a Gaussian prior.

Unlike ordinary ridge regression, Bayesian Ridge Regression does not require the regularization parameter to be specified manually. Instead, the parameters $\alpha$ and $\lambda$ are estimated from the data by maximizing the marginal likelihood.

After training, predictions for a new observation $\mathbi{x}_*$ are obtained using the posterior mean

\[
\hat{y}_*
=
\mathbi{x}_*^{\top}
\boldsymbol{\mu}.
\]

Because the model estimates an entire posterior distribution rather than a single coefficient vector, Bayesian Ridge Regression can also quantify uncertainty in both the model coefficients and the predictions.

\section*{Derivation of the Likelihood Function}

Recall that the Bayesian Ridge Regression model assumes

\[
\mathbi{y}
=
\mathbf{X}\mathbi{w}
+
\boldsymbol{\varepsilon},
\]

where

\[
\boldsymbol{\varepsilon}
\sim
\mathcal{N}
\left(
\mathbf{0},
\alpha^{-1}\mathbf{I}
\right).
\]

Rearranging the equation gives

\[
\boldsymbol{\varepsilon}
=
\mathbi{y}
-
\mathbf{X}\mathbi{w}.
\]

Since the noise vector is Gaussian and adding a constant shifts only the mean of a Gaussian random variable, it follows that

\[
\mathbi{y}
\sim
\mathcal{N}
\left(
\mathbf{X}\mathbi{w},
\alpha^{-1}\mathbf{I}
\right).
\]

Therefore, the likelihood function can be written as

\[
p(\mathbi{y}\mid\mathbf{X},\mathbi{w},\alpha)
=
\mathcal{N}
\left(
\mathbf{X}\mathbi{w},
\alpha^{-1}\mathbf{I}
\right).
\]

Using the density function of a multivariate Gaussian distribution,

\[
p(\mathbi{y}\mid\mathbf{X},\mathbi{w},\alpha)
=
(2\pi)^{-n/2}
|\alpha^{-1}\mathbf{I}|^{-1/2}
\exp
\left(
-\frac12
(\mathbi{y}-\mathbf{X}\mathbi{w})^\top
(\alpha^{-1}\mathbf{I})^{-1}
(\mathbi{y}-\mathbf{X}\mathbi{w})
\right).
\]

Since

\[
(\alpha^{-1}\mathbf{I})^{-1}
=
\alpha\mathbf{I},
\]

we obtain

\[
p(\mathbi{y}\mid\mathbf{X},\mathbi{w},\alpha)
=
(2\pi)^{-n/2}
|\alpha^{-1}\mathbf{I}|^{-1/2}
\exp
\left(
-\frac{\alpha}{2}
(\mathbi{y}-\mathbf{X}\mathbi{w})^\top
(\mathbi{y}-\mathbf{X}\mathbi{w})
\right).
\]

Ignoring constants that do not depend on $\mathbi{w}$ gives

\[
p(\mathbi{y}\mid\mathbf{X},\mathbi{w},\alpha)
\propto
\exp
\left(
-\frac{\alpha}{2}
(\mathbi{y}-\mathbf{X}\mathbi{w})^\top
(\mathbi{y}-\mathbf{X}\mathbi{w})
\right).
\]

Using the identity

\[
(\mathbi{y}-\mathbf{X}\mathbi{w})^\top
(\mathbi{y}-\mathbf{X}\mathbi{w})
=
\|\mathbf{X}\mathbi{w}-\mathbi{y}\|^2,
\]

we obtain

\[
p(\mathbi{y}\mid\mathbf{X},\mathbi{w},\alpha)
\propto
\exp
\left(
-\frac{\alpha}{2}
\|\mathbf{X}\mathbi{w}-\mathbi{y}\|^2
\right).
\]

\section*{Derivation of the Prior Distribution}

Bayesian Ridge Regression places a Gaussian prior on the coefficient vector:

\[
p(\mathbi{w}\mid\lambda)
=
\mathcal{N}
\left(
\mathbf{0},
\lambda^{-1}\mathbf{I}
\right).
\]

Using the multivariate Gaussian density,

\[
p(\mathbi{w}\mid\lambda)
=
(2\pi)^{-p/2}
|\lambda^{-1}\mathbf{I}|^{-1/2}
\exp
\left(
-\frac12
\mathbi{w}^\top
(\lambda^{-1}\mathbf{I})^{-1}
\mathbi{w}
\right).
\]

Since

\[
(\lambda^{-1}\mathbf{I})^{-1}
=
\lambda\mathbf{I},
\]

we obtain

\[
p(\mathbi{w}\mid\lambda)
=
(2\pi)^{-p/2}
|\lambda^{-1}\mathbf{I}|^{-1/2}
\exp
\left(
-\frac{\lambda}{2}
\mathbi{w}^\top\mathbi{w}
\right).
\]

Ignoring constants independent of $\mathbi{w}$ gives

\[
p(\mathbi{w}\mid\lambda)
\propto
\exp
\left(
-\frac{\lambda}{2}
\mathbi{w}^\top\mathbi{w}
\right).
\]

Because

\[
\mathbi{w}^\top\mathbi{w}
=
\|\mathbi{w}\|^2,
\]

it follows that

\[
p(\mathbi{w}\mid\lambda)
\propto
\exp
\left(
-\frac{\lambda}{2}
\|\mathbi{w}\|^2
\right).
\]

\section*{Derivation of the Posterior Distribution}

Applying Bayes' theorem gives

\[
p(\mathbi{w}\mid\mathbf{X},\mathbi{y},\alpha,\lambda)
=
\frac{
p(\mathbi{y}\mid\mathbf{X},\mathbi{w},\alpha)
p(\mathbi{w}\mid\lambda)
}{
p(\mathbi{y}\mid\mathbf{X},\alpha,\lambda)
}.
\]

Since the denominator does not depend on $\mathbi{w}$, it acts as a normalization constant and therefore

\[
p(\mathbi{w}\mid\mathbf{X},\mathbi{y},\alpha,\lambda)
\propto
p(\mathbi{y}\mid\mathbf{X},\mathbi{w},\alpha)
p(\mathbi{w}\mid\lambda).
\]

Substituting the likelihood and prior distributions,

\[
\propto
\exp
\left(
-\frac{\alpha}{2}
\|\mathbf{X}\mathbi{w}-\mathbi{y}\|^2
\right)
\exp
\left(
-\frac{\lambda}{2}
\|\mathbi{w}\|^2
\right).
\]

Using the identity

\[
e^a e^b
=
e^{a+b},
\]

we obtain

\[
p(\mathbi{w}\mid\mathbf{X},\mathbi{y},\alpha,\lambda)
\propto
\exp
\left(
-\frac{\alpha}{2}
\|\mathbf{X}\mathbi{w}-\mathbi{y}\|^2
-
\frac{\lambda}{2}
\|\mathbi{w}\|^2
\right).
\]

\section*{Computing the Posterior Mean and Covariance}

Starting from

\[
-\frac{\alpha}{2}
\|\mathbf{X}\mathbi{w}-\mathbi{y}\|^2
-
\frac{\lambda}{2}
\|\mathbi{w}\|^2,
\]

expand the squared norm:

\[
\|\mathbf{X}\mathbi{w}-\mathbi{y}\|^2
=
(\mathbf{X}\mathbi{w}-\mathbi{y})^\top
(\mathbf{X}\mathbi{w}-\mathbi{y}).
\]

Multiplying out,

\begin{align*}
(\mathbf{X}\mathbi{w}-\mathbi{y})^\top
(\mathbf{X}\mathbi{w}-\mathbi{y})
&=
(\mathbi{w}^\top\mathbf{X}^\top-\mathbi{y}^\top)
(\mathbf{X}\mathbi{w}-\mathbi{y})\\
&=
\mathbi{w}^\top\mathbf{X}^\top\mathbf{X}\mathbi{w}
-\mathbi{w}^\top\mathbf{X}^\top\mathbi{y}
-\mathbi{y}^\top\mathbf{X}\mathbi{w}
+\mathbi{y}^\top\mathbi{y}.
\end{align*}

Since

\[
\mathbi{w}^\top\mathbf{X}^\top\mathbi{y}
=
\mathbi{y}^\top\mathbf{X}\mathbi{w},
\]

the expression simplifies to

\[
(\mathbf{X}\mathbi{w}-\mathbi{y})^\top
(\mathbf{X}\mathbi{w}-\mathbi{y})
=
\mathbi{w}^\top\mathbf{X}^\top\mathbf{X}\mathbi{w}
-
2\mathbi{w}^\top\mathbf{X}^\top\mathbi{y}
+
\mathbi{y}^\top\mathbi{y}.
\]

Substituting into the exponent,

\begin{align*}
&
-\frac{\alpha}{2}
\left(
\mathbi{w}^\top\mathbf{X}^\top\mathbf{X}\mathbi{w}
-
2\mathbi{w}^\top\mathbf{X}^\top\mathbi{y}
+
\mathbi{y}^\top\mathbi{y}
\right)
-
\frac{\lambda}{2}
\mathbi{w}^\top\mathbi{w}\\
&=
-\frac12
\mathbi{w}^\top
(
\lambda\mathbf{I}
+
\alpha\mathbf{X}^\top\mathbf{X}
)
\mathbi{w}
+
\alpha
\mathbi{w}^\top
\mathbf{X}^\top
\mathbi{y}
-
\frac{\alpha}{2}
\mathbi{y}^\top\mathbi{y}.
\end{align*}

The final term does not depend on $\mathbi{w}$ and can be absorbed into the normalization constant.

Define

\[
\mathbf{\Sigma}^{-1}
=
\lambda\mathbf{I}
+
\alpha\mathbf{X}^\top\mathbf{X}.
\]

Then

\[
-\frac12
\mathbi{w}^\top
\mathbf{\Sigma}^{-1}
\mathbi{w}
+
\alpha
\mathbi{w}^\top
\mathbf{X}^\top
\mathbi{y}.
\]

Completing the square,

\[
-\frac12
(\mathbi{w}-\boldsymbol{\mu})^\top
\mathbf{\Sigma}^{-1}
(\mathbi{w}-\boldsymbol{\mu})
+
\text{constant},
\]

where

\[
\boldsymbol{\mu}
=
\alpha
\mathbf{\Sigma}
\mathbf{X}^{\top}
\mathbi{y}.
\]

Therefore,

\[
p(\mathbi{w}\mid\mathbf{X},\mathbi{y},\alpha,\lambda)
=
\mathcal{N}
(
\boldsymbol{\mu},
\mathbf{\Sigma}
),
\]

with

\[
\mathbf{\Sigma}
=
\left(
\lambda\mathbf{I}
+
\alpha\mathbf{X}^{\top}\mathbf{X}
\right)^{-1}.
\]

\section*{Relationship to Ridge Regression}

Substituting the posterior covariance into the posterior mean gives

\[
\boldsymbol{\mu}
=
\alpha
\left(
\lambda\mathbf{I}
+
\alpha\mathbf{X}^{\top}\mathbf{X}
\right)^{-1}
\mathbf{X}^{\top}
\mathbi{y}.
\]

Factor $\alpha$ from the matrix inside the inverse:

\[
\lambda\mathbf{I}
+
\alpha\mathbf{X}^{\top}\mathbf{X}
=
\alpha
\left(
\mathbf{X}^{\top}\mathbf{X}
+
\frac{\lambda}{\alpha}\mathbf{I}
\right).
\]

Using

\[
(c\mathbf{A})^{-1}
=
\frac{1}{c}
\mathbf{A}^{-1},
\]

we obtain

\[
\left(
\lambda\mathbf{I}
+
\alpha\mathbf{X}^{\top}\mathbf{X}
\right)^{-1}
=
\frac{1}{\alpha}
\left(
\mathbf{X}^{\top}\mathbf{X}
+
\frac{\lambda}{\alpha}\mathbf{I}
\right)^{-1}.
\]

Substituting into the expression for $\boldsymbol{\mu}$ gives

\begin{align*}
\boldsymbol{\mu}
&=
\alpha
\left[
\frac{1}{\alpha}
\left(
\mathbf{X}^{\top}\mathbf{X}
+
\frac{\lambda}{\alpha}\mathbf{I}
\right)^{-1}
\right]
\mathbf{X}^{\top}
\mathbi{y}\\
&=
\left(
\mathbf{X}^{\top}\mathbf{X}
+
\frac{\lambda}{\alpha}\mathbf{I}
\right)^{-1}
\mathbf{X}^{\top}
\mathbi{y}.
\end{align*}

This is identical to the ridge regression solution

\[
\mathbi{w}
=
\left(
\mathbf{X}^{\top}\mathbf{X}
+
\gamma\mathbf{I}
\right)^{-1}
\mathbf{X}^{\top}
\mathbi{y},
\]

where

\[
\gamma
=
\frac{\lambda}{\alpha}.
\]

Therefore, ridge regression can be interpreted as the posterior mean (or equivalently the MAP estimate) of Bayesian Ridge Regression under a Gaussian prior.

\end{document}